\section{Needs and Objectives}

\subsection{Needs} 
	While the distribution of information and the consistent growth in access to technology have made the feasibility and demand for rapid prototyping/consumer-aimed PCBA tools, the market for such tools is yet to be clearly defined. PCBA tools have typically been relegated to industrial use. In a commercial setting, testing physical components as opposed to simulating such components brings the benefit of not needing to model non-idealities. In a consumer setting, small-batch runs of PCBs may be unreasonable to produce due to cost. The expansion of small-scale PCBA tools can provide new avenues to both commercial and at-home users. 

\subsection{Objective}
	We aim to create a system which handles portions of the PCB manufacturing process. This machine should be affordable, precise enough to handle typical IC packages such as LQFP, and easy to maintain. The assembly of a PCB entails placing solder on the pads of an unpopulated board, placing components on their respective pads, and reflowing the solder in order to finish the process. A system to handle this process must manage the input of board information, solder paste, various ICs and board components, and the board itself in order to provide a finished assembled board. 
